\section{NDN Background and Naming Responsibilities}
\subsection{Named Retrieval, Forwarding State, and Object Identity}
NDN retrieves Data by name using Interest and Data packets~\cite{ndn,ndnpacket}. An application expresses an Interest describing the Data it seeks. A matching Data packet can be supplied by its producer or an intermediate cache. Forwarders use forwarding information to direct Interests and pending-Interest state to deliver returning Data. This permits network-layer reuse of named Data without requiring the consumer to identify its present storage location. Availability still depends on forwarding, connectivity, and a reachable copy; a name is not a delivery guarantee.

Names have several related uses in this proposal. A service name identifies the requested operation. A message name, together with the message's validated fields, identifies an invocation exchange. An object name identifies data such as an image, model fragment, or output tensor. Names also participate in trust rules that constrain acceptable signing keys~\cite{trustschema}. These uses should agree, but should not be conflated: matching a service prefix does not by itself prove service permission, and a valid producer signature does not prove that the object is the input assigned to the current execution.

For example, an illustrative image name could have the form
\begin{center}\texttt{/org/uav7/view/task42/view2/version1}\end{center}
under the sensing UAV's producer namespace. The application must define what the task, view, and version components mean and validate the signer. A processing Provider can refer to that exact image rather than rename it as its own observation. This example explains the naming principle, not a mandatory NDNSF wire format. Names and traffic patterns can remain visible even when Content is encrypted. Sensitive identifiers should therefore not be embedded in public names merely because the associated payload is protected.

The forwarding Content Store offers opportunistic caching. A repository provides an application-managed place to retain objects for later retrieval. Neither a cache hit nor a successful digest check supplies the complete admission decision for a dependent service. The consumer still needs the producer's authority, the intended object identity, and the current invocation or collaboration association. Likewise, a Data packet's FreshnessPeriod is NDN metadata that guides cache-freshness treatment; it is not a service deadline or a substitute for message expiry, current permission checks, or application replay state.

Sync helps participants learn about changes in a named dataset~\cite{li2018brief,moll2022sok}. Applications can then retrieve and validate the announced Data. Thus, publication in NDNSF means making a named object available and announcing it through the configured dissemination mechanism. It should not be read as an unsolicited push of a Data packet through arbitrary forwarders. Sync recovery can repair missing announcements or trigger retrieval while time remains. It cannot ensure completion after an application deadline or during an indefinitely disconnected interval.

These distinctions define the meaning of data-driven execution here. A validated Request, ACK, Selection, result, or declared intermediate object supplies evidence for a local state transition. Timers and local scheduling policies still exist. A timeout can close collection or terminate an operation, and a planning algorithm can decide which assignment to publish. Decisions that affect other participants are communicated through named, verifiable messages; a local callback alone cannot inform a remote participant. This does not imply that all runtime state is network Data or that packet arrival alone authorizes execution.

\subsection{From NDN Building Blocks to Service Decisions}
The Forwarding Information Base (FIB) supplies name-prefix forwarding information, the Pending Interest Table (PIT) records pending retrievals, and the Content Store (CS) holds cached Data~\cite{yi2013case}. Compatible pending Interests can share a returning Data packet. This network-level aggregation is distinct from selecting several service executors. A PIT entry also has a different lifetime from a service operation: its expiration does not cancel an application handler or erase the User's request state.

\begin{table}[htbp]\centering\small
\caption{NDN building blocks and the service decisions that remain above them.}
\label{tab:ndn-runtime-map}
\begin{tabularx}{\textwidth}{p{.23\textwidth}XX}
\toprule Building block & NDNSF use & Additional service responsibility\\\midrule
Interest/Data and names & Retrieve messages, results, and object segments & Define the operation, request identity, and completion condition\\\midrule
FIB, PIT, and CS & Forward and reuse named Data & Select work, track its lifecycle, and check acceptable object versions\\\midrule
Signatures and trust schema & Validate signing authority and Data integrity & Check service permission and assignment to this invocation\\\midrule
Dataset synchronization (generic Sync; NDN-SVS implementation) & Discover newly published message names & Validate messages and process ACK, Selection, and timeout state\\\midrule
Content encryption & Restrict access to protected Content & Define recipients, distribute keys, and apply permission changes\\\bottomrule
\end{tabularx}
\end{table}

In the software stack, the local NDN forwarder handles packet forwarding, while ndn-cxx supplies application packet, signing, validation, and retrieval interfaces~\cite{ndn-cxx}. NDN State Vector Sync (NDN-SVS) supplies a concrete dataset-synchronization implementation~\cite{ndnsvs-github}. NDNSF organizes these facilities into service interactions. Its service policies do not require forwarders to interpret GPU load, flight readiness, or a collaboration role. Those meanings belong to the framework and application contracts.

\section{Remote Procedure Call (RPC), Robotics Middleware, and Messaging Alternatives}
\subsection{RPC Interfaces and Microservice Integration}
Remote procedure call (RPC) frameworks provide typed interfaces for invoking operations on another process or host. gRPC is one RPC implementation that supports unary and streaming calls~\cite{grpc}. Robotics middleware such as ROS provides reusable communication abstractions for robot components~\cite{ros,ros2}, while message-oriented middleware supports decoupled producers and consumers~\cite{messagequeue}. These systems are relevant alternatives to application-specific networking. NDNSF does not claim that endpoint discovery, asynchronous interaction, or distributed application composition is absent from existing middleware.

Endpoint resolution identifies a destination, but does not necessarily establish that the destination is authorized and currently willing to execute a particular invocation. Middleware or the application may add those policies; this proposal treats their expression and evidence as part of the comparison.

The comparison concerns how the same application requirements are expressed. An RPC experiment must state how endpoints are resolved, connections are managed, retries are performed, and failures are detected in the selected implementation. An NDN experiment must state how names are discovered, how configured synchronization and retransmission operate, how Data is validated, and how Content Store behavior is handled. A typed call is not inherently less secure than named-object exchange; security depends on the accompanying authentication, authorization, and encryption mechanisms. Likewise, network-layer caching is useful only when object reuse is permitted by the workload and operation semantics. Reusing an image does not mean that an invocation with physical side effects may be repeated safely.

For NDNSF, the relevant question is how service decisions and application data can remain linked through names and verifiable objects. The evaluation must test that link under stated workloads rather than equate a difference in networking abstraction with superiority. The named-computation and command-authorization systems discussed below provide more direct protocol precedents.

\subsection{Robot Components and Asynchronous Messaging}
Robot software separates camera, telemetry, control, and perception components, often combining service calls with continuously updated publish--subscribe topics~\cite{ros,ros2}. NDNSF's service-container design adopts this separation between component behavior and communication support. It does not replace the flight controller or imply that robotics middleware cannot implement authorization. The comparison instead asks how an incoming operation identifies the intended vehicle and how stored observations remain associated with their producer and operation.

Message brokers provide asynchronous delivery and buffering for logs, telemetry, and workflow events~\cite{messagequeue}. The relevant design questions are where data is retained, who authorizes publication and consumption, and how delivery acknowledgment differs from task completion. A broker may be appropriate for a stable deployment. NDN's named retrieval provides a different placement and validation model, but still requires an explicit retention policy when late retrieval matters. Neither messaging abstraction alone chooses a willing execution Provider or defines safe repetition of a physical command.

\subsection{Microservices and Edge Computing}
CFEC studies computation offloading across vehicular fog, edge, and cloud resources using mobility and resource information~\cite{cfec}. MIA-NDN studies microservice-aware Interest aggregation and PIT lifetime management based on input characteristics~\cite{miandn}. SECaaS studies the deployment of security functions at the multi-access edge, including access-control enforcement and anonymization~\cite{secaas}. These works widen the comparison beyond function invocation: resource placement, forwarding state, and deployable security functions are relevant parts of the surrounding system.

Differences in these systems' objectives do not, by themselves, establish unique NDNSF capabilities. Placement feasibility does not establish invocation permission, parameter-aware aggregation does not determine the assigned predecessor in a collaboration, and deploying an access-control service is a different question from authorizing calls to that service. Conversely, NDNSF's application-layer checks do not replace these systems' resource-management or forwarding mechanisms. A claim of superiority would require a common workload and matching requirements; their published results are not interchangeable baselines for the experiments in this proposal.
